% GENERATED FILE. Edit canonical lessons, then run npm run generate:books.
% canonical-source-hash: fnv1a64:ab7d73a3a9089102
% canonical-lessons: TE-C530-viyyankudu, TE-C530-muttata, TE-C530-menalludu, TE-C530-cuttam, TE-C530-avva

\chapter{In-Laws and Kin}
\label{ch:te-in-laws-and-kin}
\hlchaptermodality{530}

\begin{chapteropening}
\textbf{By the end of this chapter:} \emph{I can name a co-father-in-law, a great-grandfather, a nephew and my relations.}

Complete the last lesson of chapter 530: I can name a co-father-in-law, a great-grandfather, a nephew and my relations.
\end{chapteropening}

\section[\texorpdfstring{\te{వియ్యంకుడు} (viyyaṅkuḍu)}{వియ్యంకుడు (viyyaṅkuḍu)}]{\te{వియ్యంకుడు} (viyyaṅkuḍu) — a co-father-in-law; one's child's father-in-law}

\label{lesson:TE-C530-viyyankudu}



\begin{quote}
Before the new one: say the Telugu for a mother's elder sister, then the Telugu for a niece.
\end{quote}

\subsection*{You'll want to know: \te{వియ్యంకుడు}}
\textbf{\te{వియ్యంకుడు}} — \emph{viyyaṅkuḍu} — \textquotedblleft{}a co-father-in-law; one's child's father-in-law\textquotedblright{}.

\begin{quote}
\te{పెళ్ళి} \te{తరువాత} \te{నాన్న}, \te{వియ్యంకుడు} \te{మంచి} \te{స్నేహితులు} \te{అయ్యారు}.
\end{quote}

\begin{quote}
\emph{peḷḷi taruvāta nānna, viyyaṅkuḍu mañci snēhitulu ayyāru.}
\end{quote}

\begin{quote}
After the wedding, dad and his co-father-in-law became good friends.
\end{quote}

The fathers of a married couple are each other's \te{వియ్యంకుడు}; the mothers are \te{వియ్యపురాలు}. \te{వియ్యం} is the bond between the two families.

\begin{cousinweb}
\textbf{\te{వియ్యంకుడు}} (\emph{viyyaṅkuḍu}) is listed in DEDR's appendix of Indo-Aryan loans (DEDR App. 58) beside \emph{viyyamu}, \textquotedblleft{}a marriage alliance\textquotedblright{}; the family goes back to Sanskrit \emph{vivāha}, \textquotedblleft{}marriage\textquotedblright{} (CDIAL 11920).
\end{cousinweb}

\begin{grammarlens}[title={What you've built}]
One more word for describing what you lived and felt.
\end{grammarlens}

\subsection*{Your turn}
\begin{itemize}
\raggedright
  \item \emph{Say it:} \emph{viyyaṅkuḍu}
  \item \emph{Say it:} \emph{viyyaṅkuḍu}, once more
  \item \emph{Say it:} \emph{peḷḷi taruvāta nānna, viyyaṅkuḍu mañci snēhitulu ayyāru.} — After the wedding, dad and his co-father-in-law became good friends.
  \item \emph{Recall:} say the Telugu for a mother's elder sister, then the Telugu for a niece, then say \emph{viyyaṅkuḍu} again
\end{itemize}

\begin{tcolorbox}[breakable,colback=teal!4,colframe=teal!35!black,title={Before you move on}]
What does \te{వియ్యంకుడు} mean? (\textquotedblleft{}A co-father-in-law; one's child's father-in-law\textquotedblright{}.) Say it once more.
\end{tcolorbox}
\section[\texorpdfstring{\te{ముత్తాత} (muttāta)}{ముత్తాత (muttāta)}]{\te{ముత్తాత} (muttāta) — a great-grandfather}

\label{lesson:TE-C530-muttata}



\begin{quote}
Before the new one: say the Telugu for a niece, then the Telugu for a co-father-in-law.
\end{quote}

\subsection*{You'll want to know: \te{ముత్తాత}}
\textbf{\te{ముత్తాత}} — \emph{muttāta} — \textquotedblleft{}a great-grandfather\textquotedblright{}.

\begin{quote}
\te{మా} \te{ముత్తాత} \te{కట్టిన} \te{ఇల్లు} \te{ఇంకా} \te{ఉంది}.
\end{quote}

\begin{quote}
\emph{mā muttāta kaṭṭina illu iṅkā undi.}
\end{quote}

\begin{quote}
The house my great-grandfather built is still standing.
\end{quote}

The father of a grandfather; a great-grandmother is \te{ముత్తవ్వ}.

\begin{cousinweb}
\textbf{\te{ముత్తాత}} (\emph{muttāta}) is native Dravidian: DEDR 4954 lists Telugu \emph{muttāta}, \textquotedblleft{}great-grandfather\textquotedblright{}, and \emph{mutt-avva}, \textquotedblleft{}great-grandmother\textquotedblright{}, in the article of Tamil \emph{mutu}, \textquotedblleft{}old\textquotedblright{}: the old grandfather.
\end{cousinweb}

\begin{grammarlens}[title={What you've built}]
One more word for describing what you lived and felt.
\end{grammarlens}

\subsection*{Your turn}
\begin{itemize}
\raggedright
  \item \emph{Say it:} \emph{muttāta}
  \item \emph{Say it:} \emph{muttāta}, once more
  \item \emph{Say it:} \emph{mā muttāta kaṭṭina illu iṅkā undi.} — The house my great-grandfather built is still standing.
  \item \emph{Recall:} say the Telugu for a niece, then the Telugu for a co-father-in-law, then say \emph{muttāta} again
\end{itemize}

\begin{tcolorbox}[breakable,colback=teal!4,colframe=teal!35!black,title={Before you move on}]
What does \te{ముత్తాత} mean? (\textquotedblleft{}A great-grandfather\textquotedblright{}.) Say it once more.
\end{tcolorbox}
\section[\texorpdfstring{\te{మేనల్లుడు} (mēnalluḍu)}{మేనల్లుడు (mēnalluḍu)}]{\te{మేనల్లుడు} (mēnalluḍu) — a nephew; a sister's son}

\label{lesson:TE-C530-menalludu}



\begin{quote}
Before the new one: say the Telugu for a co-father-in-law, then the Telugu for a great-grandfather.
\end{quote}

\subsection*{You'll want to know: \te{మేనల్లుడు}}
\textbf{\te{మేనల్లుడు}} — \emph{mēnalluḍu} — \textquotedblleft{}a nephew; a sister's son\textquotedblright{}.

\begin{quote}
\te{నా} \te{మేనల్లుడు} \te{కాలేజీలో} \te{చదువుతున్నాడు}.
\end{quote}

\begin{quote}
\emph{nā mēnalluḍu kālējīlō caduvutunnāḍu.}
\end{quote}

\begin{quote}
My nephew is studying at college.
\end{quote}

A man's sister's son or a woman's brother's son; \te{అల్లుడు} alone is a son-in-law.

\begin{cousinweb}
\textbf{\te{మేనల్లుడు}} (\emph{mēnalluḍu}) joins native \emph{alluḍu}, \textquotedblleft{}son-in-law\textquotedblright{} (DEDR 2410), to \emph{mēna-}, from DEDR's appendix of Indo-Aryan loans (DEDR App. 53).
\end{cousinweb}

\begin{grammarlens}[title={What you've built}]
One more word for describing what you lived and felt.
\end{grammarlens}

\subsection*{Your turn}
\begin{itemize}
\raggedright
  \item \emph{Say it:} \emph{mēnalluḍu}
  \item \emph{Say it:} \emph{mēnalluḍu}, once more
  \item \emph{Say it:} \emph{nā mēnalluḍu kālējīlō caduvutunnāḍu.} — My nephew is studying at college.
  \item \emph{Recall:} say the Telugu for a co-father-in-law, then the Telugu for a great-grandfather, then say \emph{mēnalluḍu} again
\end{itemize}

\begin{tcolorbox}[breakable,colback=teal!4,colframe=teal!35!black,title={Before you move on}]
What does \te{మేనల్లుడు} mean? (\textquotedblleft{}A nephew; a sister's son\textquotedblright{}.) Say it once more.
\end{tcolorbox}
\section[\texorpdfstring{\te{చుట్టం} (cuṭṭaṁ)}{చుట్టం (cuṭṭaṁ)}]{\te{చుట్టం} (cuṭṭaṁ) — a kinsman; a relation}

\label{lesson:TE-C530-cuttam}



\begin{quote}
Before the new one: say the Telugu for a great-grandfather, then the Telugu for a nephew.
\end{quote}

\subsection*{You'll want to know: \te{చుట్టం}}
\textbf{\te{చుట్టం}} — \emph{cuṭṭaṁ} — \textquotedblleft{}a kinsman; a relation\textquotedblright{}.

\begin{quote}
\te{పెళ్ళికి} \te{చుట్టాలంతా} \te{వచ్చారు}.
\end{quote}

\begin{quote}
\emph{peḷḷiki cuṭṭālantā vaccāru.}
\end{quote}

\begin{quote}
All the relations came to the wedding.
\end{quote}

Plural: \te{చుట్టాలు}, \emph{cuṭṭālu}.

Plural \te{చుట్టాలు}; \te{చుట్టరికం}, kinship.

\begin{cousinweb}
\textbf{\te{చుట్టం}} (\emph{cuṭṭaṁ}) is native Dravidian: DEDR 2715 lists Telugu \emph{cuṭṭamu}, \textquotedblleft{}relation, kinsman\textquotedblright{}, beside \emph{cuṭṭu}, \textquotedblleft{}to go round, to encircle\textquotedblright{}: the people around you.
\end{cousinweb}

\begin{grammarlens}[title={What you've built}]
One more word for describing what you lived and felt.
\end{grammarlens}

\subsection*{Your turn}
\begin{itemize}
\raggedright
  \item \emph{Say it:} \emph{cuṭṭaṁ}
  \item \emph{Say it:} \emph{cuṭṭaṁ}, once more
  \item \emph{Say it:} \emph{peḷḷiki cuṭṭālantā vaccāru.} — All the relations came to the wedding.
  \item \emph{Recall:} say the Telugu for a great-grandfather, then the Telugu for a nephew, then say \emph{cuṭṭaṁ} again
\end{itemize}

\begin{tcolorbox}[breakable,colback=teal!4,colframe=teal!35!black,title={Before you move on}]
What does \te{చుట్టం} mean? (\textquotedblleft{}A kinsman; a relation\textquotedblright{}.) Say it once more.
\end{tcolorbox}
\section[\texorpdfstring{\te{అవ్వ} (avva)}{అవ్వ (avva)}]{\te{అవ్వ} (avva) — a granny; an old woman}

\label{lesson:TE-C530-avva}



\begin{quote}
Before the new one: say the Telugu for a nephew, then the Telugu for a kinsman.
\end{quote}

\subsection*{You'll want to know: \te{అవ్వ}}
\textbf{\te{అవ్వ}} — \emph{avva} — \textquotedblleft{}a granny; an old woman\textquotedblright{}.

\begin{quote}
\te{పల్లెలో} \te{మా} \te{అవ్వ} \te{ఒంటరిగా} \te{ఉంటుంది}.
\end{quote}

\begin{quote}
\emph{pallelō mā avva oṇṭarigā uṇṭundi.}
\end{quote}

\begin{quote}
Our granny lives alone in the village.
\end{quote}

Rural and affectionate; also any old woman.

\begin{cousinweb}
\textbf{\te{అవ్వ}} (\emph{avva}) is native Dravidian: DEDR 273 lists Telugu \emph{avva}, \textquotedblleft{}mother, grandmother, an old woman\textquotedblright{}, with Tamil \emph{avvai}.
\end{cousinweb}

\begin{grammarlens}[title={What you've built}]
One more word for describing what you lived and felt.
\end{grammarlens}

\subsection*{Your turn}
\begin{itemize}
\raggedright
  \item \emph{Say it:} \emph{avva}
  \item \emph{Say it:} \emph{avva}, once more
  \item \emph{Say it:} \emph{pallelō mā avva oṇṭarigā uṇṭundi.} — Our granny lives alone in the village.
  \item \emph{Recall:} say the Telugu for a nephew, then the Telugu for a kinsman, then say \emph{avva} again
\end{itemize}

\begin{tcolorbox}[breakable,colback=teal!4,colframe=teal!35!black,title={Before you move on}]
What does \te{అవ్వ} mean? (\textquotedblleft{}A granny; an old woman\textquotedblright{}.) Say it once more.
\end{tcolorbox}
